\begin{tikzpicture}[
  font=\scriptsize,
  panah/.style={-{Stealth[length=2mm]}, thick},
  tahap/.style={draw, rounded corners, fill=blue!8, text width=1.45cm,
    minimum height=1.05cm, align=center, inner sep=2pt},
  hasil/.style={tahap, fill=green!10},
  gerbang/.style={draw, rounded corners, fill=orange!12, font=\tiny,
    align=center, inner sep=2pt},
  pita/.style={draw, rounded corners=2pt, fill=praditagreen!15},
  ket/.style={font=\tiny, gray!60!black}
]
  % Enam tahap pipeline dari commit sampai produksi.
  \node[tahap] (t1) at (0,0) {\textit{Commit} dan tinjauan kode};
  \node[tahap] (t2) at (1.9,0) {\textit{Build} dan uji otomatis};
  \node[tahap] (t3) at (3.8,0) {Pemindaian keamanan};
  \node[tahap] (t4) at (5.7,0) {Artefak bertanda versi};
  \node[tahap] (t5) at (7.6,0) {\textit{Staging} dan uji penerimaan};
  \node[hasil] (t6) at (9.5,0) {Produksi bertahap};
  \foreach \asal/\tujuan in {t1/t2, t2/t3, t3/t4, t4/t5, t5/t6} {
    \draw[panah] (\asal) -- (\tujuan);
  }

  % Umpan balik cepat ke pengembang.
  \draw[panah, gray!70, dashed] (t3.north) to[bend right=28]
    node[ket, above] {umpan balik dalam hitungan menit} (t1.north);

  % Gerbang keputusan rilis pada continuous delivery.
  \node[gerbang] (gerbang) at (8.55,1.15) {keputusan rilis\\(manusia)};
  \draw[thick, orange!70!black] (gerbang.south) -- (8.55,0.08);

  % Cakupan tiga istilah.
  \draw[pita] (-0.8,-1.08) rectangle (6.5,-0.68);
  \node[font=\tiny] at (2.85,-0.88)
    {\textit{Continuous Integration}: setiap \textit{commit} dibangun dan diuji};
  \draw[pita] (-0.8,-1.63) rectangle (8.4,-1.23);
  \node[font=\tiny] at (3.8,-1.43)
    {\textit{Continuous Delivery}: selalu siap rilis, manusia memutuskan kapan};
  \draw[pita] (-0.8,-2.18) rectangle (10.3,-1.78);
  \node[font=\tiny] at (4.75,-1.98)
    {\textit{Continuous Deployment}: lolos semua tahap, langsung ke produksi};
\end{tikzpicture}
